\section{Related work}\label{sec:related-work}

The mathematical background of this paper lies first in the classical study of Cantor and fractal objects generated by fixed symbolic, self-similar, or graph-directed rules. Standard references include Falconer’s general account of fractal geometry and dimension theory, the original IFS and Moran-construction sources, and the graph-directed and symbolic framework developed by Mauldin and Urba\'nski \cite{Hutchinson1981Fractals,Moran1946AdditiveFunctions,FalconerFractalGeometry,MauldinUrbanskiGDMS}. In those settings, the central object is typically specified in advance by a fixed rule system, and the main questions concern geometry, dimension, pressure, and invariant measures for that fixed object \cite{Bowen1979Quasicircles,Ruelle1978ThermodynamicFormalism,WaltersErgodic1982}.

A second background line is thermodynamic formalism beyond strictly additive observables. The pressure side of the present paper is much closer in spirit to nonadditive or almost-subadditive thermodynamic constructions than to a fixed self-similar pressure formula, even though our final object is not presented in that language alone. For general background on nonadditive thermodynamic formalism, see \cite{BarreiraNonadditiveTF,CaoFengHuang2008SubAdditive,Mummert2006AlmostAdditive,Kingman1968Subadditive}. The present paper does not attempt to prove a broad nonadditive formalism for arbitrary evolving Cantor substrates; instead it closes a pressure object on a specific audited shell-stable class.

A third line of context is the study of local or context-dependent iterated constructions. Recent work on local iterated function systems provides a useful point of comparison because it highlights how admissible compositions may depend on local geometry rather than a single global IFS rule \cite{OliveiraVarandasLocalIFS,BarnsleyDemkoEltonGeronimo1988,FalconerRandomFractals1986,RempeGillenUrbanski2016NonAutonomousCIFS}. However, the theorem package closed in this paper is not a broad local-IFS theorem. We do not claim a shell-general local-IFS theorem or a general non-SFT breadth theorem. Our claims are restricted to the audited shell on which the canonical hybrid object is closed.

Within the prior six-birds program, three earlier papers provide the relevant background and vocabulary. The first is the foundational closure/saturation/extension formalism of \cite{TsiokosFoundations}. The second is the substrate-generation perspective developed in \cite{TsiokosCreateStone}, which motivates the use of a continuous six-mechanism substrate rather than a fixed discrete pipeline. The third is the closure-deficit / conditional-disintegration perspective in \cite{TsiokosCastStone}, which is the closest conceptual antecedent of the thermodynamic consequence theoremlet proved here. These papers are cited only for background and historical context. The present paper is not a general six-birds framework paper; it is a Cantor mathematics paper built on a specific canonical hybrid object.

The continuous six-mechanism substrate machinery itself is adapted from prior six-birds/PICA-style work and is not the locus of novelty. What is new in this branch is the closed theorem package supported by that substrate on the audited shell: a continuous full-loop lawfulness theoremlet, a cocycle pressure closure theoremlet, a strict theory-extension theoremlet, and a conditional pressure disintegration theoremlet on a canonical hybrid Cantor object.

Two routes are explicitly outside the present package. First, no broader-class theorem beyond the audited shell is claimed, even though several wider routes were explored during development. Second, no direct stratumwise root-separation theorem is claimed; the thermodynamic consequence is formulated as a conditional disintegration of the base pressure over packaging fibers \cite{Rokhlin1949FundamentalIdeas,LedrappierWalters1977RelativisedVP,CoverThomas2006InfoTheory}.
